\chapter{About why must exist the bijection $\Omega$}

	\noindent
	Imagine we define A, as the set of all possible t-uplas, with finite size, of Natural numbers.
	\\\\
	
	\noindent
	It is a known result: A is countable. You could have your own bijections, I am able to build my own CLJA for that set. It takes around 90 pages :D, so let's say it is a known result, okey? It is true for all the mathematical community: you can do it with potences of prime numbers (if I am not wrong, I use to confuse easily bijections and injections, they have the same use for me).
	\\\\
	
	\noindent
	If A is countable:\\
	A X A  is countable.\\
	A X A X A X A is countable too.\\
	The union (A X A) U (A X A X A X A) is countable too.\\
	Think calmly: LCF is a subset of that last union. LCF MUST BE countable.
	\\\\
	
	\noindent
	Taking ONLY the ordered versions of those t-uplas, is a subset of A.\\
	$\{ (0), (1) \}$, the DR values, are a subset of A... a very small one, but a subset of A.
	\\\\
	
	\noindent
	The bijection:\\
	$\Omega: (LCF \:\: U \:\: Previous) \rightarrow \mathbb{N}$\\
	MUST EXISTS !!
	\\\\
	
	\noindent
	In the unnecesary comment 0010, you have the formula of $\Omega$ written in spanish (sorry).
	\\\\
	
	\noindent
	Here you have the python code that implements direct and inverse functions ( it needs a lot of refactoring, and my energy is used to create one version after another of THIS DOCUMENT :D, and others...):\\
	https://github.com/CLJAs/clja-ftc
	\\\\
	
	\noindent
	And here you have the first 300 millions cases of $\Omega$, listed. Each zip contains 100 million cases, each text file contains one million cases(lines). We have used the inverse function to calculate the member of LCF from a natural number. And after that, we use the direct function, to calculate again the natural number, FROM THAT member of LCF. The initial and final natural number must be the same. OK mark means it is:\\
	https://drive.google.com/file/d/1kcEdhZn-UwNz5cLEZI9QyhqfzAV4mim3/view?usp=sharing \\
	https://drive.google.com/file/d/1EfJib9Vokcz6CAU4UBJxYfC9y1oZQMsJ/view?usp=sharing \\
	https://drive.google.com/file/d/1nAXuF6oK42TXiprn-nEDLB36Q56dMP2Z/view?usp=sharing \\
	\\\\
	
	\noindent
	Like $\Omega$ defines a bijection, between $LCF$ and $\mathbb{N}$, we can "translate" any possible natural number, to members of LCF, defining exactly the same partitions and sub-partitions inside $\mathbb{N}$. We will not worry about paths between $LCF$ and $\mathbb{N}$. It's enough just with reaching subsets of $LCF$. Our paths of the Divide and Conquer strategy, can end in $LCF$.
	\\\\ 
	
	
	
	
	
	
	
	
	 

\newpage